\begin{lemma}[Singletons belong to a nontrivial full-base tree]
Let $F$ be a front on $\mathbb N$ and suppose $\emptyset\notin F$.  Then
every singleton $\{n\}$ belongs to $T(F)$.
\end{lemma}
\begin{proof}
The assumption $\emptyset\notin F$ excludes the trivial alternative in the
base clause of \noderef{Front}; hence $\mathsf{FrontBase}(F)=\mathbb N$.
Fix $n\in\mathbb N$ and consider the infinite tail
$Y=\{n\}\cup\{k\mid n<k\}$.  It is a subset of the full base.  Density in
\noderef{Front} gives $t\in F$ with $\mathsf{ProperPrefixSet}(t,Y)$, using
the paper notation $t\segs Y$ from \noderef{ProperPrefixSet}.  The witness
$t$ is nonempty: otherwise $t=\emptyset$, contradicting $t\in F$ and the
hypothesis $\emptyset\notin F$.

Unfolding \noderef{ProperPrefixSet}, choose $m\in Y$ such that

\[
t=\{k\in Y\mid k<m\}.
\]

The number $n$ is the least element of $Y$.  If $m=n$, the displayed set is
empty, contrary to the nonemptiness of $t$; hence $n<m$, and consequently
$n\in t$.  We now verify explicitly that $\{n\}$ is an initial segment of
$t$ in the sense of \noderef{InitialSegment}.  If $t=\{n\}$, this is the
equality case.  Otherwise, the finite set $t\setminus\{n\}$ is nonempty;
let $q$ be its least element.  Every element of $t$ other than $n$ is at
least $q$, while $n<q$, so
\[
\{k\in t\mid k<q\}=\{n\}.
\]
Thus $\{n\}$ is an initial segment of $t$.  The definition of
\noderef{PrefixTree} now gives $\{n\}\in T(F)$.
\end{proof}
